\chapter{Background}%
\label{ch:background}

% In this chapter, you provide information that has a benefit for an informed audience.
% Often, you use this to explain a technique or a concept you use, referring to the literature.
% This is different from Related Work (\cref{ch:related}), where you show the relation of literature to your contributions (to ground, compare, or inform).
% A good rule of thumb is that your thesis should still make sense if the background chapter is removed.
This chapter introduces the concepts that form the foundation of this thesis. It covers the vulnerability management context in which the work is situated, the architectural principles that motivate the design approach, the main ideas behind explainability, and the human factors that make cybersecurity communication difficult in practice. The chapter does not yet position this thesis against prior solutions in detail; that is the role of Chapter~\ref{ch:related}. Instead, the goal here is to clarify the concepts and terminology needed to understand the problem and the proposed approach.

\section{Vulnerability Management as a Communication Problem}

Vulnerability management is a recurring cybersecurity process in which organisations identify, assess, prioritise, and remediate weaknesses in their systems. In practice, this process is often supported by scanners and management platforms such as Qualys, which generate structured findings containing information such as CVE identifiers, severity scores, affected assets, exploitability indicators, and remediation status. These findings are useful because they provide a relatively standardised technical basis for planning mitigation and tracking exposure over time.

However, the output of a vulnerability scanner is not a decision in itself. It is an input to a larger process in which risks must be prioritised, remediation must be scheduled, trade-offs must be justified, and findings must often be communicated to people who do not work directly with the tooling. In a consulting or enterprise context, those stakeholders may include delivery managers, sales staff, project leads, product owners, compliance stakeholders, and clients. Their concerns are typically not phrased in terms of raw vulnerability metadata, but in terms of questions such as: What is the actual risk? Why is this important now? What should be done next? What happens if it is delayed? Which evidence supports this recommendation?

This creates a translation problem. The scanner produces technically precise output, but the decision context requires representations that are more selective, contextual, and understandable for the intended audience. The problem becomes harder when the same technical findings need to be communicated to different stakeholder groups with different goals. An engineer may want technical detail and evidence. A manager may want prioritisation and resource implications. A client-facing stakeholder may need a concise, defensible explanation that connects technical findings to business impact.

Research on intrusion detection and security operations shows that even technical practitioners face substantial effort when interpreting and configuring security tooling \cite{werlinger2008ids}. This is relevant here because it suggests that difficulties in interpreting cybersecurity output do not disappear simply by giving stakeholders more data or more dashboards. The challenge is not only the amount of information, but also how it is structured, framed, and connected to the decisions people must make.

\section{Software Architecture, Abstraction, and Stakeholder Views}

Software architecture provides a useful lens for understanding this problem. A software architecture defines the high-level structure of a system, including its components, responsibilities, and interactions. Beyond technical structure, architecture also serves as a mechanism for managing complexity, reasoning about quality attributes, and communicating design decisions to different stakeholders \cite{bass2012software}.

A central architectural principle is abstraction. Rather than exposing every implementation detail, an architecture defines boundaries and representations that make a system understandable at different levels of detail. Parnas argued that modularisation is effective when modules are organised around design decisions that should be hidden from other parts of the system \cite{parnas1972modules}. This idea remains important in modern systems because it supports changeability, maintainability, and conceptual clarity. In the context of this thesis, the same principle can be applied to cybersecurity communication: raw technical findings should not be exposed directly to every stakeholder in the same form. Instead, the system should provide abstractions that present only the level of detail needed for a particular audience and purpose.

Architecture also addresses the fact that different stakeholders need different views of the same system. Kruchten’s 4+1 view model illustrates that architectural descriptions should not assume one universal representation, but rather a set of complementary views that address different concerns \cite{kruchten1995viewmodel}. Although this model was developed in a general software design context, the underlying idea is directly relevant here. Vulnerability information can be seen as one underlying technical reality that needs to be represented differently depending on whether the audience is a security analyst, a delivery manager, or a non-technical client stakeholder.

Bass, Clements, and Kazman also emphasise that architectural design should be driven by quality attributes such as modifiability, usability, and performance \cite{bass2012software}. For this thesis, several of these quality attributes are especially relevant. \textit{Traceability} matters because explanations should remain connected to the underlying scanner evidence. \textit{Modifiability} matters because cybersecurity tooling and reporting needs change over time. \textit{Usability} matters because the explanations must support real interpretation rather than merely rephrasing technical outputs. \textit{Maintainability} matters because explanation logic that is hard-coded in ad hoc reports or dashboards quickly becomes inconsistent and expensive to evolve.

Taken together, these architectural ideas support the core design assumption of this thesis: explainability in cybersecurity communication should not be treated only as a reporting feature or interface layer, but as an architectural concern that shapes how data is transformed, represented, and delivered across the system.

\section{Explainability and Interpretability}

Explainability and interpretability are often used interchangeably, but the literature shows that the concepts are broader and less straightforward than they first appear. Doshi-Velez and Kim argue that interpretability should not be treated as a vague intrinsic property of a model, but should instead be evaluated in relation to a specific user and task \cite{doshivelez2017rigorous}. Liao et al.\ complement this by showing empirically that users of AI systems have diverse and task-specific explanation needs, and that these needs must be elicited rather than assumed \cite{liao2020questioning}. Lipton similarly criticises broad claims about interpretability and distinguishes between different goals, such as transparency, post-hoc explanation, and trust \cite{lipton2018mythos}. This is important for this thesis because it avoids treating explainability as something that is automatically achieved by adding a specific technique.

Survey work in XAI shows that a wide range of explanation methods exist, including feature attribution, surrogate models, rule extraction, example-based explanations, and counterfactual explanations \cite{guidotti2018survey, barredo2020xai}. Popular examples include LIME, which approximates local model behaviour with an interpretable surrogate \cite{ribeiro2016should}; SHAP, which attributes output contributions to features using game-theoretic reasoning \cite{lundberg2017unified}; and counterfactual explanations, which show how an input would need to change to obtain a different outcome \cite{wachter2017counterfactual}. These methods have contributed significantly to the practical XAI landscape, but they are mostly concerned with explaining predictions made by machine learning models. Their limitations and implications for the problem studied in this thesis are discussed in more detail in Chapter~\ref{ch:related}.

For this thesis, two observations from the explanation literature are especially relevant. First, explanations are not only technical objects; they are also communicative acts addressed to an audience. Miller argues that human explanations are typically contrastive and selective \cite{miller2019explanation}. People usually do not want an exhaustive account of all causes. Instead, they want an explanation of why one outcome occurred rather than another, and they want the most relevant factors, not the entire mechanism. In a vulnerability management setting, this means that a useful explanation is unlikely to be a full dump of every scanner field. More often, it is a concise rationale that answers a decision-oriented question such as why one vulnerability is prioritised over another or why a patch should be treated as urgent.

Second, explanation can be approached as a system-level design concern rather than only as a model-level feature. Barredo Arrieta et al.\ note that explainability depends on the audience and the purpose of the explanation, not just on the computational method used to produce it \cite{barredo2020xai}. This aligns closely with the goals of this thesis. The objective is not to make an underlying detection model inherently interpretable, but to design software abstractions that organise technical evidence into stakeholder-oriented explanations.

\section{Data-to-Text and Explanation as Translation}

A useful complement to XAI comes from natural language generation and data-to-text research. This literature studies how structured data can be transformed into natural language summaries and narratives that are understandable to human users. Reiter and Dale describe this as a process that includes content selection, aggregation, and linguistic realisation \cite{reiter2000nlg}. Gatt and Krahmer show that these techniques have been applied in domains such as weather reporting, healthcare, and business analytics, where raw structured data needs to be communicated in a form that supports understanding and action \cite{gatt2018survey}.

This perspective matters because it treats explanation not only as the exposure of internal reasoning, but also as a translation process. In other words, explanation can involve deciding which technical details matter for a given user, how they should be grouped into higher-level claims, and how they should be phrased so that the result is understandable without losing its connection to the source evidence.

For cybersecurity communication, this suggests a shift in focus. Instead of only asking how to explain a classifier, it becomes possible to ask how a system can transform technical findings into explanations that are stakeholder-appropriate, selective, and traceable. In this thesis, this idea is important because the proposed artifact includes an explanation layer that does more than expose low-level output. It also performs abstraction and translation.

\section{Human Factors and Usable Security}

Cybersecurity systems operate in socio-technical environments, which means that their effectiveness depends not only on detection quality but also on how people interpret and act on their outputs. Cranor’s framework for reasoning about the human in the loop highlights that security outcomes depend strongly on human perception, attention, and decision-making \cite{cranor2008framework}. This is highly relevant for this thesis because vulnerability reporting often fails not because the underlying scanner is wrong, but because the resulting information does not fit the stakeholder’s mental model, role, or decision context.

Usable security research repeatedly shows that security mechanisms can create extra burden when they are technically correct but poorly aligned with the people expected to use them. In practice, technical precision alone is not enough. Information must also be understandable, appropriately framed, and presented in a way that supports action. For non-technical stakeholders, raw security indicators such as CVSS scores, exploit flags, and vulnerability identifiers often require additional interpretation before they become meaningful. For technical stakeholders, too much abstraction can be problematic if it removes needed evidence or obscures important uncertainty.

This creates a design tension. Explanations should simplify, but not oversimplify. They should support understanding, but remain tied to evidence. They should be role-specific, but consistent enough to be maintained across the system. These tensions motivate the architectural approach taken in this thesis.

\section{Design Science Research}

The research in this thesis follows a Design Science Research approach. Design Science Research is appropriate when the goal is to create and evaluate an artifact intended to solve a relevant problem \cite{hevner2004dsr}. In Software Engineering, this is a common and legitimate research format because many contributions are not only empirical observations but also new designs, methods, architectures, or prototypes intended to improve practice.

The artifact in this thesis is a conceptual software architecture and prototype for transforming vulnerability scanner outputs into stakeholder-oriented explanation objects. The focus is therefore not on proving a universal theory of explainability, but on designing and evaluating a concrete solution in context. This fits well with the Design Science perspective because the research outcome is both constructive and evaluative: a design is proposed, justified, implemented in prototype form, and assessed against stakeholder needs and quality concerns.

In summary, the background of this thesis combines four ideas. Vulnerability management creates a communication problem because technical outputs must support decisions made by different stakeholders. Software architecture provides principles for managing this problem through abstraction, modularisation, and stakeholder views. Explainability literature clarifies that useful explanations are audience-dependent, selective, and often contrastive. Finally, usable security and data-to-text research show that explanation must be designed with human interpretation in mind. These ideas together motivate the design-oriented approach taken in the remainder of the thesis.